% ============================================================
%  H. Brøchner, "Problemet om Tro og Viden" (1868)
%  Forord (pp. 1–7)
%  TRANSCRIPTION (Danish) — from OCR of KB scan
% ============================================================
\documentclass[12pt,a4paper]{article}
\usepackage[utf8]{inputenc}
\usepackage[T1]{fontenc}
\usepackage[danish]{babel}
\usepackage{microtype}
\usepackage{setspace}
\usepackage[margin=1.2in]{geometry}
\usepackage{fancyhdr}
\usepackage{hyperref}

\hypersetup{
  pdftitle={Problemet om Tro og Viden — Forord},
  pdfauthor={H. Brøchner},
  hidelinks
}

\pagestyle{fancy}
\fancyhf{}
\rhead{Brøchner, \emph{Problemet om Tro og Viden} (1868)}
\cfoot{\thepage}

\setstretch{1.3}

\title{\textbf{Problemet om Tro og Viden}\\[6pt]
  \large En historisk-kritisk Afhandling\\[10pt]
  \normalsize Forord}
\author{H.\ Brøchner\\[4pt]
  \small Professor i Philosophie}
\date{Kjøbenhavn: P.\ G.\ Philipsens Forlag, 1868}

\begin{document}
\maketitle
\thispagestyle{fancy}

\bigskip

\noindent Dette Skrift gjengiver, i alt Væsentligt uforandret,
Indholdet af de Forelæsninger, jeg i Efteraarshalvaaret 1867 har holdt
ved Kjøbenhavns Universitet over Forholdet mellem Tro og Viden, mellem
Religion, Philosophie og Ethik. Disse Forelæsninger omfattede
Undersøgelsens historisk-kritiske Deel og dannede Forberedelssen til et
Forsøg paa en selvstændig Løsning af det behandlede Problem, der i
Løbet af dette Aar vil blive offentliggjort.

Den historisk-kritiske Deel af Undersøgelsen giver, foruden en kort
Skizze af Problemets Tilbliven, en udførlig Kritik af de
Løsningsforsøg, der hos os have faaet en særegen Betydning, og af
hvilke eet --- det, der hidrører fra S.\ Kierkegaard --- er gjort med
en Originalitet, en Alvor og Dybde, der til alle Tider vil bevare det
denne Betydning. Gjennem Kritiken af disse Løsningsforsøg, som under
deres Udvikling og Begrundelse have staaet i et vedvarende polemisk
eller mæglende Forhold til de tidligere og andetsteds fremkomne, og
deels kunne betragtes som Resultater af tidligere udprægede
Hovedretninger, deels væsentligen have fuldbyrdet Kritiken af
saadanne, ere de betydeligste af de ældre og samtidige Løsningsforsøg
belyste og tildels de fuldstændige Forudsætninger givne for en endelig
Dom om dem; og denne historiske Kritik vil saa finde sit Supplement i
det Skrift, der knytter sig til dette.

Problemets nærværende Stilling er for vor Literaturs Vedkommende
betegnet ved Kampen imellem en Form af Theologien, der væsentligen
laaner de Vaaben, med hvilke den vil hævde Foreningen af Tro og Viden,
paa eklektisk Maade fra nyere philosophiske Blandingsformer, og det
antitheologisk-religiøse Standpunkt, der ved den principielle
Modsætning mellem Tro og Viden vil sikre Religionens og Videnskabens
Gyldighed i deres gjensidige Uafhængighed og Muligheden til deres
fredelige Samværen i Bevidstheden, medens den i Kraft af samme
principielle Modsætning frafjender Theologien al Ret til at bære
Videnskabens Navn, hvilket den tilegner sig idet den gjør Fordring paa
at være en ,,Troesvidenskab''.

Den philosopherende Theologie, der fastholder Muligheden af
Troesindeholdets Optagen og Udprægen i Videns Form, og derfor ikke
blot vil hævde Theologien Videnskabens Navn, men tillægger den, som
Aabenbaringsvidenskabens afsluttende Form, Principatet blandt
Videnskaberne og Dommermyndigheden over dem, har til sit historiske
Udgangspunkt havt det speculative System, der gjennem en tvetydig
Udvidelse af ,,Religionens'' Begreb bragte Religionen i dens specifike
Betydning i et identisk Forhold til Philosophien, fra hvilken ikke en
Indholdsforskjel, men kun en Formsdifferents skulde sondre den.
Efterat denne tvetydige Enhed har maattet opgives under de Kampe, der
udviklede sig efter den Hegelske Religionsphilosophies Fremkomst, har
den af Philosophien paavirkte Theologie, idet den ombyttede
Speculationens Betegnelse med det Christeliges, dog bevaret
Mæglingstendentsen; men den støtter sig nu nærmest paa en anden
philosophisk Retning, og særlig paa det System, der ved at bringe en
Fordobling ind i Philosophien skaffer Plads for en ,,fri Tænkning'',
der fra sin nære Slægtning: Phantasien, laaner den ,,Fjederham'',
hvilken den kan overflyve Mysteriets Afgrunde; og den anbringer
forsigtigt et Forbehold, der afskjærer ogsaa den ,,højere Videnskabs''
Konseqvenser, naar Tanken, trods ,,Friheden'', uvilkaarlig kommer til
at følge sine væsensnødvendige Love.

Den antitheologiske Theorie, der ligeledes har den Hegelske Philosophie
til Forudsætning, men hvis Forhold til Forudsætningen oprindelig er
negativt, polemisk, havde til historisk Incitament ogsaa den Protest,
der fra en Philosophie, som paa een Gang optræder som Konseqvents af
og som Modsætning til hiin speculative Mæglingslære, nedlagdes mod
Enheden af Religion og Philosophie. Som den betydeligste Repræsentant
for denne Retning staaer Feuerbach --- en Mand, hvis Skrifter kun
meget ufuldstændigt ere kjendte og kun meget overfladisk ere bedømte
af dem, der hos os have kritiseret ham. Idet han opfattede Religionens
Standpunkt som praktisk i Modsætning til Videnskabens theoretiske, og
idet han nærmere bestemte hint ved en bestandig skarpere Udhæven af
det alogiske Moment deri, fik han den uforenelige Modsætning til
Synspunkt for Forholdet og fordrede i den sande Erkjendelses og den
sande ethiske Handlings Interesse Religionens Absorption i Lidenskab
og Ethik. Religionen som saadan skulde, som et Alogisk, som et
ufuldkomment Gjennemgangstrin i Udviklingen af Menneskeaandens
Selvbevidsthed, forsvinde; kun dens Navn blev tilbage i uegentlig og
overført Betydning.

Enig med Feuerbach i Modsætningen, men uenig i Begrundelsen og
Anvendelsen af denne, gjør saa S.\ Kierkegaard Sondringen mellem
Religion og Philosophie i Religionens Interesse. Menneskeaanden, der i
Existentsens Uro ikke kan gribe det Evige i Erkjendelsens objective
Vished, og selv erkjender sin Erkjendelses Skranke, finder kun i det
religiøse Forhold, i Troen, en subjectiv existentiel Vished om det,
der er dens dybeste Interesse; den finder først der sin sande
Forsoning. Videns Gyldighed benægtes ikke, men den objective Viden
begrændses til et Omraade, paa hvilket Aandens dybeste Interesser ikke
finde deres Tilfredsstillelse. Tro og Viden faaer saaledes hver sit
Gebeet, og derved og ved Videns Subordination skal Forenligheden af
dem i samme Bevidsthed blive mulig.

Denne Kierkegaardske Opfattelse af Problemet om Forholdet mellem Tro
og Viden er bleven optagen af Professor R.\ Nielsen og fremsat paa en
saadan Maade, at man har villet stille ham i Originalitet lige med
Theoriens Ophavsmand og har tillagt ham Løsningen af en Opgave af
samme Oprindelighed som den, hiin stræbte at løse. Hvor lidet
berettiget denne Anskuelse er, vil i dette Skrift blive paaviist. Det
vil vise sig, at Problemet ved Prof.\ Nielsen ikke er blevet ført
videre, at der ved den Form, han har givet Løsningen, kun er bragt
Usikkerhed og Vaklen ind i de afgjørende Bestemmelser, at Forsøget paa
en psychologisk og metaphysisk Begrundelse af Troesprincipets Sondring
fra Vidensprincipet som absolut uensartet med dette er mislykket og
fører til uløselige Modsigelser, og at det Samme gjælder om hans
Begrundelse af den Sætning, at de absolut uensartede Principier skulle
kunne forenes i samme Bevidsthed.

Naar der da, ved Siden af S.\ Kierkegaard, den egentlige Repræsentant
for det antitheologisk-religiøse Standpunkt, bliver taget saa udførligt
Hensyn til Prof.\ Nielsen, saa er det ikke fordi Theorien hos ham har
faaet en dybere Begrundelse eller en rigere Udvikling, men fordi den
ved ham har ombyttet den esoteriske Læres Charakter med den exoteriske,
er trængt ind i ,,de Dannedes'' vide Kreds, og i denne fra en
Vidensag er ifærd med at blive til en Troessag og at fastslaaes som
Dogme i Kraft af et inappellabelt $\alpha\dot{\upsilon}\tau\grave{o}\varsigma\ \check{\varepsilon}\varphi\eta$.

Som Repræsentant for den philosopherende Theologie fremtræder i dette
Skrift Biskop Martensen; men idet Udviklingen nøie knytter sig til det
i hans Skrifter Givne, er ingensteds det Almene, det Typiske i
Standpunktet tabt af Sigt, eller glemt over det Individuelle. Til andre
theologiske Forsøg hos os har det ikke været nødvendigt directe at
tage Hensyn.

Den endelige Opgave, hvis Løsning forberedes gjennem Kritiken, der
umiddelbart kun kan give et negativt Resultat, er Hævdelsen af et
Standpunkt, der i Forhandlingerne hos os ikke er kommet til Orde, og
som kun i Former, der ikke kunne tilfredsstille, er gjort gjældende
andetsteds. Dette Standpunkt fastholder en saadan Modsætning mellem
Videnskaben, særlig Philosophien, og de positive Religionsformer,
mellem Viden og Tro i Betydning af Aabenbaringstro, at Viden og denne
Tro ikke kunne bringes til Enhed i samme Bevidsthed; --- det fastholder
Umuligheden af, at en sand og indholdsfyldig, det hele, virkelige
Menneskeliv omfattende Ethik kan udfolde sig paa Basis af den positive
Religion med dens supranaturale Handlingslov; --- og det fastholder den
indre, uopløselige Enhed mellem den sande Erkjendelse og den sande
ethiske Handlen, dennes nødvendige Betingethed ved hiin. Men medens
det fastholder hiin Modsætning, hævder det paa den anden Side ligesas
afgjørende, idet det sondrer det Religiøses Væsensbestemmelse fra de
positive Religioners phænomenale Former, i hvilke det kun seer et
partielt og ufuldkomment, endeliggjørende Udtryk for hiin, idet det
altsaa sondrer det i Religionerne Religiøse fra Religionerne i deres
historiske Former, den indre Enhed af det Religiøse, det Philosophiske
og det Ethiske, og gjennem denne indre Enhed Menneskelivets indre
Enhed, Betingelsen for dets sande og fulde Forsoning, Forsoningen
indenfor Virkeligheden, den ene og udelte Mennesknaturs Forsoning med
sig i sin Væsensgrund. Idet det hævder Væsentligheden og
Nødvendigheden af det religiøse Forhold for Mennesket, hvis
Selvbevidsthed først bliver fuldstændig virkeliggjort i dette, gjennem
Gudsbevidtstheden, og som ifølge sit Væsens Grundforhold aldrig har
været og aldrig vil være uden Religion; --- idet det tillige seer den
ethiske Handlen og den sande Erkjendelse som vindende en høiere
Betydning gjennem Enheden med det religiøse Grundforhold; --- og idet
det lader den sande Viden og den sande ethiske Handlen blive en
adæqvat Virkelighedsform for det Religiøse, opfatter det denne Viden
og denne Handlen som i sig indesluttende et Høieste og Ubetinget, i
hvilket det i Religionerne Væsentlige er bevaret. Dette Standpunkt
betegne vi som \textit{den religiøst forsonede humane Bevidsthed}.

Det adskiller sig fra det af Feuerbach indtagne ikke blot derved, at
dets metaphysiske Grundbegreb er et andet --- et saadant, paa hvilket
en Erkjendelse og en Handlen med Idealitetens og det Almengyldiges
Præg kan bygges, og ud fra hvilket Begrebet om den menneskelige Frihed
kan hævdes og et religiøst Forhold begrundes til det Absolute, med
hvilket Mennesket er i Væsensenhed --- men ogsaa derved, at det
kæmper mod Theologien og Dogmet ikke blot i Ethikens og Videnskabens,
men først og fremmest i selve Religionens, i det religiøse Forholds
Uendeligheds Interesse.

Det adskiller sig fra den gamle Hegelske Skoles Opfattelse af
Religionens Forhold til Erkjendelsen, idet det ikke seer det Religiøse
blot som en Videns-Form, men som Udtrykket for et det hele menneskelige
Væsen i dets Concretion omfattende Grundforhold, og idet det ikke
lader Aandslivets Medium være et illusorisk abstract Evighedselement,
men Virkeligheden med dens reale Betingelser og Skranker.

Med den antitheologisk-religiøse Theorie mødes det i Opfattelsen af
Theologiens Forhold til Videnskaben. Idet det hviler paa en
videnskabelig Erkjenden og i det udelte Menneskevæsen giver Erkjendelsen
Principatet, maa det være dets Opgave at hævde Videnskabens klare
Begreb og dens i dens eget Væsen begrundede Autonomie mod
Tvetydigheder og paa Tvetydigheder støttede Forsøg paa at underlægge
Videnskaben under en fremmed, uensartet Autoritets Herredømme. Men
medens der saaledes i et væsentligt Punkt er en Overensstemmelse
mellem det og hiin Theorie, træder det i en afgjort Opposition til
denne, naar Forholdet mellem Tro og Viden, mellem Religion, Philosophie
og Ethik skal bestemmes. Denne Opposition gjælder Kierkegaards
Opfattelse, med Hensyn til hvilken det vil blive viist, at den rokker
al Erkjendelses Vished og at den lader det Ethiske forsvinde i et
abstract og virkelighedsløst Religiøst. Men den gjælder endnu
bestemtere den Form, i hvilken Theorien er bleven udført og fremstillet
af Nielsen --- en Form, i hvilken Kierkegaard sikkerlig ikke vilde
vedkjende sig den. Den Form, i hvilken Nielsen udpræger Theorien,
fører, trods alle de mere eller mindre behændige Vendinger, der gjøres,
og ved hvilke denne Konseqvents søges undgaaet, til en factisk Dualisme
i Menneskevæsenet, til en Deling af Aandens Liv fra Grunden af ved de
,,absolut uensartede'' Principier, til Bevidsthedens uundgaaelige
Selvopløsning; den forvanker det Ethiske i dets Grund og aabner
Adgangen til det Religiøse for alt det Superstitionøse og Absurde. Med
den Overbevisning, at den af Nielsen statuerede absolute Uensartethed
mellem Videns Princip paa den ene Side, Troens og Handlingens paa den
anden er ,,en for Religiøsiteten, Sædeligheden og al sand
Karakterudvikling fordærvelig Misforstaaeelse'', har jeg underkastet
hans Theorie en gjennemgaaende Kritik, og søgt at godtgjøre dens
Uholdbarhed, dens mangelfulde Begrundelse, dens Selvmodsigelser, dens
fordærvelige Konseqvenser.

Idet jeg offentliggjør dette Bidrag til Undersøgelsen af et Problem,
der gjennem Aarhundreder har havt og bestandig vil have den meest
afgjørende Betydning for Menneskeaanden, er det mit Haab, at det maa
give Anledning til en fornyet, alvorlig og samvittighedsfuld Drøftelse,
til en skarp og aaben Kamp mellem Principier. En saadan Kamp, selv om
den føres med skarpe Vaaben, kan kun være i Sagens Interesse, og man
therefore være kjær for den Kæmpende, for hvem Sagen, ikke et
egoistisk Selvhensyn, er det Høieste. En fast og dyb personlig
Overbevisning skal jeg altid agte; ved gode Grunde skal jeg altid være
ligesaa paavirrkelig, som jeg vil være upaavirrkelig ved Andet end
disse.

\end{document}
